% Nouveau EGA V (cote 25): archive cote 25, site page 144 (PDF page 145).
% typescript. Draft compared: none.
% AI-transcribed by Claude Opus 5.5 (claude-opus-5-5, adaptive thinking, effort high), Message Batches API, batch ega25b (msgbatch_01A7n5YuXubUPNK1zsepi8RG), 2026-11-13; not yet checked by eye.
\scanpage{25/144}
\mspage{\hand{5}}

La colonne provient de la suite exacte de transtitivité pour les
morphismes lisses $X\to Y$ et $Y\to S$ , en \var{\xout{xx}}{}{typewriter strike-out}{ega25_25_144-v0001} tensorisant par $\underline{O}_{V1}$ (cela reste
exact, car tous les modules de la suite sont loc libres). La ligne
horizontale est un\var{\xout{e}}{}{typewriter strike-out}{ega25_25_144-v0002} cas particulier d'une suite exacte obtenue chaque
fois que sur X \var{\deleted{\var{\xout{lisse}}{}{typewriter strike-out}{ega25_25_144-v0004}}}{}{deletion}{ega25_25_144-v0003} sur S, on a une section \unsure{$\phi$} d'un module loc libre F
et qu'on prend le schéma W des zéros; on trouve une suite exacten
\[ \overset{\hand{\vee}}{F}\otimes \underline{O}_X \longrightarrow \underline{\Omega}^1_{X/S}\otimes \underline{O}_W \longrightarrow \underline{\Omega}^1_{W/S} \longrightarrow 0 \]
et \var{\inserted{(si X/S est lisse)}}{(si X/S est lisse)}{insertion}{ega25_25_144-v0005} le premier homomorphisme est injectif exactement en les points où
W est lisse sur X avec la ``bonne '' \unsure{dimension} relative (i.e. partout dans
\var{\inserted[typed]{\var{\xout{pratiquement}}{}{typewriter strike-out}{ega25_25_144-v0007}}}{\var{\xout{pratiquement}}{}{typewriter strike-out}{ega25_25_144-v0007}}{insertion}{ega25_25_144-v0006}
le cas actuel). \var{\xout{Cet homomorphisme}}{}{typewriter strike-out}{ega25_25_144-v0008} Cette suite exacte est \var{\xout{écrite en}}{}{typewriter strike-out}{ega25_25_144-v0009}
\var{\xout{par 16 dans le cas où F = $\underline{O}_X$, la généralisation}}{}{typewriter strike-out}{ega25_25_144-v0010} \var{\xout{(à ceci près}}{}{typewriter strike-out}{ega25_25_144-v0011} une
conséquence immédiate de la suite exacte
\[ \underline{J}/\underline{J}^2\longrightarrow \underline{\Omega}^1_{X/S}\otimes \underline{O}_W \longrightarrow \underline{\Omega}^1_{X/S} \longrightarrow 0 \]
\var{\margin{16 \textbar}}{\placedmargin{16 \textbar}}{marginal note placed}{ega25_25_144-v0012} qui figure au par.16 (on y pourrait mettre en corollaire la version
signalée ici). La caractérisation de l'ensemble des points où on
peut mettre un zéro à gauche est contenue dans le critère \var{\over{j}{J}}{J}{writing over}{ega25_25_144-v0013}acobien.

Notons que l'on a un isomorphisme canonique
\[ \underline{P}^1_{X/S} = \underline{\Omega}^1_{X/S} + \underline{O}_X \ ,\ \text{d'où}\ \check{\underline{P}}^1_{X/S} = \hand{\unsure{\mathcal{G}|\mathcal{T}}}_{X/S} + \underline{O}_X \ , \]
d'autre part on \unsure{vérifi\var{\over{e}{é}}{é}{writing over}{ega25_25_144-v0014}} que l'homomorphisme \var{\inserted{composé}}{composé}{insertion}{ega25_25_144-v0015} \hand{$\mu$} du diagramme \hand{\unsure{(1)}} est nul
sur le facteur $\underline{O}_{V1}$ , et sur le facteur $\hand{\unsure{\mathcal{G}|\mathcal{T}}}_{X/S}\otimes\underline{O}_{V1}$ \var{\inserted{il}}{il}{insertion}{ega25_25_144-v0016} se réduit à
l'homomorphisme \var{\inserted{$M(\phi)'$}}{$M(\phi)'$}{insertion}{ega25_25_144-v0017} déduit de la section $M(\phi)$ de $\underline{\mathrm{Sym}}^2(\hand{\underline{\Omega}}^1_{X/S})\otimes \underline{O}_{V1}$ déjà
envisagée. Donc en un point x de X, $M(\phi)$ est non dégénéré i.e. $M(\phi)'$
est surjectif sss \hand{$\mu$} est surjectif \var{\inserted{en x}}{en x}{insertion}{ega25_25_144-v0018}, et on voit sur le diagramme \hand{(1)},
que cela équivaut aussi à dire que \unsure{\hand{$\nu$}} est surjectif \var{\inserted{en x}}{en x}{insertion}{ega25_25_144-v0019} (car l'un et l'autre
signifie\var{\inserted{nt}}{nt}{insertion}{ega25_25_144-v0020} que l'homomorphisme canonique de\var{\xout{xx deux}}{}{typewriter strike-out}{ega25_25_144-v0021} la somme des deux sous
Modules envisagés de $\underline{\Omega}^1_{X/S}\otimes \underline{O}_{V1}$ dans ce derniers est surjectif en
x). On trouve ainsi :

\var{\margin{144}}{\placedmargin{144}}{marginal note placed}{ega25_25_144-v0022}
